\appendix
\setcounter{equation}{0}
\renewcommand{\theequation}{A.\arabic{equation}}
\renewcommand{\theHequation}{A.\arabic{equation}}

\section{Proofs and Boundary Cases}
\label{app:proofs}

\begin{proof}[Proof of Lemma~\ref{lem:demand}]
Write \(\gamma=1-\alpha-\beta>0\). Maximizing log utility in
\eqref{eq:outer-problem} gives
\begin{equation}
 \frac{\alpha}{M}=\lambda P_M,\qquad
 \frac{\beta}{N}=\lambda,\qquad
 \frac{\gamma}{Z}=\lambda.
 \label{eq:outer-focs}
\end{equation}
The budget constraint implies \(\lambda=1/Y\), hence
\begin{equation}
 M=\frac{\alpha Y}{P_M},\qquad N=\beta Y,\qquad Z=\gamma Y.
 \label{eq:outer-demands}
\end{equation}
Strictly positive shares, prices, and income make this solution interior.
Substitution yields
\begin{equation}
 U=\alpha^\alpha\beta^\beta\gamma^\gamma
 \frac{Y}{P_M^\alpha},
 \label{eq:indirect-utility-proof}
\end{equation}
which establishes \eqref{eq:real-income}.

For the inner problem, let \(\mathcal C(M_A,M_P)\) denote the right side of
\eqref{eq:channel-ces}. A corner cannot minimize cost. For example, as
\(M_A\downarrow0\) with \(M_P>0\),
\(\mathcal C_A/\mathcal C_P\to\infty\); moving along the isoquant has marginal
cost \(p_A-p_P\mathcal C_A/\mathcal C_P<0\) near the boundary. The symmetric
argument excludes \(M_P=0\). We may therefore use the interior Lagrangian
\begin{equation}
\mathcal L=p_AM_A+p_PM_P+\mu[M-\mathcal C(M_A,M_P)].
 \label{eq:channel-lagrangian}
\end{equation}
Its first-order conditions imply
\begin{equation}
 \frac{M_A}{M_P}
 =\frac{1-a}{a}\left(\frac{p_A}{p_P}\right)^{-\eta}.
 \label{eq:relative-channel-demand}
\end{equation}
Using the binding constraint gives
\begin{equation}
 M_A=(1-a)\left(\frac{p_A}{P_M}\right)^{-\eta}M,\qquad
 M_P=a\left(\frac{p_P}{P_M}\right)^{-\eta}M,
 \label{eq:conditional-demands-proof}
\end{equation}
and unit cost is \eqref{eq:price-index}. Strict quasiconcavity of the CES
aggregator makes its expenditure-requirement set strictly convex, so the
positive-price solution is unique. All terms in \eqref{eq:platform-share} are
positive, hence \(0<s<1\) for finite \(z\).

Let \(Q=(1-a)p_A^{1-\eta}+ap_P^{1-\eta}\). Since
\(d\ln p_P/dz=-1\),
\begin{equation}
 \frac{d\ln P_M}{dz}
 =\frac{ap_P^{1-\eta}}{Q}\frac{d\ln p_P}{dz}=-s.
 \label{eq:price-derivative-proof}
\end{equation}
Also,
\begin{equation}
 \frac{s}{1-s}=\frac{a}{1-a}
 \left(\frac{p_P}{p_A}\right)^{1-\eta}
 \quad\Longrightarrow\quad
 s'=(\eta-1)s(1-s).
 \label{eq:share-derivative-proof}
\end{equation}
Differentiating \eqref{eq:platform-retention} and \eqref{eq:retention} gives
\eqref{eq:retention-derivatives}.
\end{proof}

\begin{proof}[Proof of Lemma~\ref{lem:equilibrium}]
For fixed \(z\), \(T\) maps \(\mathbb R_+\) into itself because
\(B+\omega G>0\). Since \(s\in(0,1)\) and both channel-retention shares lie in
\([0,1]\), \(0\leq\omega\leq1\). For any \(Y,\widetilde Y\geq0\),
\begin{equation}
 |T(Y)-T(\widetilde Y)|
 =(\beta+\alpha\omega)|Y-\widetilde Y|,
 \qquad 0\leq\beta+\alpha\omega\leq\alpha+\beta<1.
 \label{eq:contraction-proof}
\end{equation}
The contraction theorem gives a unique fixed point and global convergence of
\(Y_{n+1}=T(Y_n)\). Solving the fixed-point equation gives
\eqref{eq:income-solution}; its numerator and denominator are positive.
Lemma~\ref{lem:demand} then uniquely determines all remaining objects.
\end{proof}

\begin{proof}[Proof of Proposition~\ref{prop:multiplier}]
Fixing \(z\), \(\kappa\), and \(G\), treat realized \(B\) as an independent
autonomous-income injection, equivalently a change in the level shifter
\(\bar B\). Differentiating \eqref{eq:income-solution} in that injection and in
gross external retail spending \(G\) gives
\eqref{eq:local-income-multiplier} and \eqref{eq:recursive-multiplier}; the
series converges because
\(\beta+\alpha\omega<1\). Moreover,
\begin{equation}
 m_0\omega-\mathcal M_G(\omega)=
 \frac{\alpha\omega(1-\omega)}
 {(1-\alpha-\beta)(1-\beta-\alpha\omega)}.
 \label{eq:multiplier-gap}
\end{equation}
Every denominator is positive, proving the inequality and both equality cases.
Direct differentiation gives \eqref{eq:retention-elasticity} and
\begin{equation}
 \mathcal M_G''(\omega)=
 \frac{2\alpha(1-\beta)}
 {(1-\beta-\alpha\omega)^3}>0.
 \label{eq:multiplier-convexity}
\end{equation}
\end{proof}

\begin{proof}[Proof of Proposition~\ref{prop:thresholds}]
At fixed \(\kappa\), \(\Delta(\kappa)>0\) and
Lemma~\ref{lem:demand} imply
\(\omega_1<\omega_0\), \(D_1>D_0\), and \(P_{M1}<P_{M0}\). With \(G=0\),
\begin{equation}
 \frac{Y_1}{Y_0}=\frac{B_1/D_1}{B_0/D_0}
 =\frac{\mathcal B}{\mathcal D},\qquad
 \frac{R_1}{R_0}=
 \frac{\mathcal B}{\mathcal D\mathcal P^\alpha}.
 \label{eq:threshold-proof-ratios}
\end{equation}
Because \(0<\mathcal P^\alpha<1\),
\(\mathcal D\mathcal P^\alpha<\mathcal D\). Comparing
\(\mathcal B\) with these ordered thresholds proves the three regions and both
boundaries. If \(Y_1/Y_0>1\), then
\(R_1/R_0=(Y_1/Y_0)/\mathcal P^\alpha>1\), excluding the fourth sign pattern.
Finally, \(\mathcal B=e^{\varepsilon(z_1-z_0)}\geq1\). Holding the other
ratios fixed and allowing \(\varepsilon\geq0\), a \(\mathcal B\) in the
joint-loss interval exists exactly when
\(\mathcal D\mathcal P^\alpha>1\).
\end{proof}

\begin{proof}[Proof of Corollary~\ref{cor:marginal}]
For \(G=0\), \(Y=B/D\), \(D'=-\alpha\omega'\), and
\begin{equation}
 \frac{d\ln Y}{dz}=\varepsilon-\frac{D'}D
 =\varepsilon-\Lambda(s,\kappa),\qquad
 \frac{d\ln R}{dz}=\frac{d\ln Y}{dz}+\alpha s.
 \label{eq:marginal-proof}
\end{equation}
Let \(d_A=1-\beta-\alpha\rho_A\geq1-\alpha-\beta>0\). Since
\begin{equation}
 D=d_A+\alpha s\Delta(\kappa),\quad
 \Delta_\kappa=-\delta_\rho,\quad
 D_\kappa=-\alpha s\delta_\rho,
 \label{eq:capacity-denominator}
\end{equation}
direct differentiation gives \eqref{eq:capacity-income} and
\eqref{eq:capacity-drag}. From
\(\mathcal M_G=\omega/D\), \(\omega_\kappa=s\delta_\rho\), and
\(D+\alpha\omega=1-\beta\), we obtain
\eqref{eq:capacity-multiplier}.

Write \(h=(\eta-1)(1-s)>0\). Then
\begin{equation}
 \Lambda(s,\kappa)
 =\frac{\alpha s h\,\Delta(\kappa)}
 {d_A+\alpha s\Delta(\kappa)}.
 \label{eq:capacity-drag-rewrite}
\end{equation}
Equation~\eqref{eq:capacity-drag} makes this function continuous and strictly
decreasing in \(\kappa\). Hence the endpoint condition around any \(q\)
implies a unique solution in \((0,1)\). Moreover,
\(\Lambda(s,0)>q\geq0\) implies \(q<h\), and solving
\(\Lambda=q\) gives
\[
 \Delta(\kappa)=\frac{q\,d_A}{\alpha s(h-q)}.
\]
Using
\(\Delta(\kappa)=\rho_A-\underline\rho_P-\kappa\delta_\rho\)
gives \eqref{eq:capacity-threshold}.

The function \(\kappa(q)\) is strictly decreasing in \(q\). Therefore the
combined endpoint condition gives
\(\kappa_R^*<\kappa_Y^*\). Equation~\eqref{eq:marginal-effects} and strict
monotonicity of \(\Lambda\) give all three sign regions and the equality
boundaries. At \(\kappa_R^*\), real income is locally constant while nominal
income falls; at \(\kappa_Y^*\), nominal income is locally constant while real
income rises.
\end{proof}

\begin{proof}[Proof of Corollary~\ref{cor:peak}]
Fix \(\kappa\) with \(\Delta(\kappa)>0\), and write
\[
 d_A=1-\beta-\alpha\rho_A>0,\qquad
 b=\alpha\Delta(\kappa)>0,\qquad
 C=\alpha(\eta-1)\Delta(\kappa)>0.
\]
Because \(D=d_A+bs\), direct differentiation gives
\begin{align}
 \Lambda_s(s,\kappa)
 &=C\frac{d_A(1-2s)-bs^2}{(d_A+bs)^2},
 \label{eq:drag-first-derivative}\\
 \Lambda_{ss}(s,\kappa)
 &=-\frac{2C\,d_A(d_A+b)}{(d_A+bs)^3}<0.
 \label{eq:drag-second-derivative}
\end{align}
Thus \(\Lambda\) is strictly concave. The numerator
\(H(s)=d_A(1-2s)-bs^2\) of \eqref{eq:drag-first-derivative} satisfies
\(H(0)>0\), \(H(1)<0\), and \(H'(s)=-2d_A-2bs<0\). It therefore has one
zero in \((0,1)\). Solving \(H(s)=0\) yields
\eqref{eq:penetration-peak}; the strict inequality follows because
\(\sqrt{1+b/d_A}>1\).

Let \(t=\sqrt{1+b/d_A}>1\). At \(s^\dagger=1/(1+t)\), use
\(b=d_A(t^2-1)\) to obtain \eqref{eq:maximum-drag}. Finally,
\(\Delta_\kappa=-\delta_\rho\) implies
\begin{align}
 \frac{\partial s^\dagger}{\partial\kappa}
 &=\frac{\alpha\delta_\rho}{2d_A t(1+t)^2}>0,
 \label{eq:peak-location-capacity}\\
 \frac{\partial\Lambda_{\max}}{\partial\kappa}
 &=-\frac{(\eta-1)\alpha\delta_\rho}{d_A t(1+t)^2}<0.
 \label{eq:peak-height-capacity}
\end{align}
These derivatives apply while \(\Delta(\kappa)>0\).
\end{proof}

\subsection{Boundary and limiting cases}

For finite \(z\), Lemma~\ref{lem:demand} rules out channel corners. Instead,
\(z\to-\infty\) implies \(s\to0\) and \(\omega\to\rho_A\), while
\(z\to\infty\) implies \(s\to1\) and
\(\omega\to\rho_P(\kappa)\).
Likewise, \(a\to0\) and \(a\to1\) produce the respective degenerate-share
limits, outside Assumption~\ref{ass:primitives}.
Because \(s\in(0,1)\), \(\omega=0\) exactly when
\(\rho_A=\rho_P(\kappa)=0\), giving \(\mathcal M_G=0\) and
\(Y=B/(1-\beta)\). Likewise, \(\omega=1\) exactly when
\(\rho_A=\rho_P(\kappa)=1\), giving \(\mathcal M_G=m_0\).

The endpoints \(\kappa=0,1\) produce
\(\rho_P(0)=\underline\rho_P\) and
\(\rho_P(1)=\overline\rho_P\), respectively. In the excluded
limit \(\delta_\rho=0\), organizational capacity has no effect. If
\(\rho_A=\rho_P(\kappa)\), retention is locally invariant to \(z\). If
\(\rho_A<\rho_P(\kappa)\), then \(\omega_z>0\); with \(G=0\) and
\(\varepsilon\geq0\), nominal and real income rise. As \(\alpha\to0\),
\[
 Y\to\frac{B+\omega G}{1-\beta},\qquad
 \mathcal M_G\to\frac{\omega}{1-\beta};
\]
as \(\beta\to0\), \(D\to1-\alpha\omega\) and
\(\mathcal M_G\to\omega/(1-\alpha\omega)\). Setting \(\varepsilon=0\) isolates
retention and price effects. If \(\alpha+\beta\uparrow1\), contraction ceases
to be uniform; at equality a finite fixed point remains only when
\(\omega<1\), while \(D=0\) at full retention.

As \(\eta\to1\),
\begin{equation}
 P_M\to p_A^{1-a}p_P^a,\qquad s\to a,
 \label{eq:cobb-douglas-limit}
\end{equation}
the Cobb--Douglas channel limit. As \(\eta\to\infty\) at fixed unequal prices,
demand converges to the cheaper channel. At equal prices the
perfect-substitutes problem has multiple
cost-minimizing allocations, although the CES sequence selects \(s=a\).

For Corollary~\ref{cor:peak}, \(\Delta(\kappa)\downarrow0\) implies
\(s^\dagger\uparrow1/2\) and \(\Lambda_{\max}\downarrow0\). At
\(\Delta(\kappa)=0\), \(\Lambda\equiv0\), so every share maximizes the zero
function and the peak is not unique. If \(\Delta(\kappa)<0\), integration
raises rather than lowers local retention; the drag-peak result does not apply.

Holding nominal \(G>0\) fixed, the general derivatives are
\begin{equation}
 \frac{d\ln Y}{dz}
 =\frac{\varepsilon B+\omega_zG}{B+\omega G}
 +\frac{\alpha\omega_z}{1-\beta-\alpha\omega}.
\label{eq:general-g-derivative}
\end{equation}
\begin{equation}
 \frac{d\ln R}{dz}
 =\frac{\varepsilon B+\omega_zG}{B+\omega G}
 +\frac{\alpha\omega_z}{1-\beta-\alpha\omega}+\alpha s.
 \label{eq:general-g-real-derivative}
\end{equation}
For organizational capacity,
\begin{equation}
 \frac{\partial\ln Y}{\partial\kappa}
 =s\delta_\rho\left[
 \frac{G}{B+\omega G}+\frac{\alpha}{D}\right]>0,\qquad
 \frac{\partial\ln R}{\partial\kappa}
 =\frac{\partial\ln Y}{\partial\kappa}.
 \label{eq:general-capacity-derivative}
\end{equation}
For a finite change define
\begin{equation}
 \mathcal N=
 \frac{B_1+\omega_1G}{B_0+\omega_0G}.
 \label{eq:general-finite-numerator}
\end{equation}
Then \(Y_1/Y_0=\mathcal N/\mathcal D\) and
\(R_1/R_0=\mathcal N/(\mathcal D\mathcal P^\alpha)\); replacing
\(\mathcal B\) by \(\mathcal N\) gives the general thresholds.
